\documentclass[11pt]{ctexart}

\usepackage[margin=2.5cm]{geometry}
\usepackage{graphicx}
\usepackage{booktabs}
\usepackage{amsmath}
\usepackage[hidelinks]{hyperref}

\setmainfont{Times New Roman}
\setCJKmainfont{SimSun}[BoldFont=SimHei, ItalicFont=KaiTi, BoldItalicFont=SimHei]

\newcommand{\figdir}{figures}

\title{\bfseries 用于 AI 生成图像检测的可解释四分支注意力融合网络}
\date{}

\begin{document}
\maketitle

\begin{abstract}
文本到图像生成模型的快速发展，使得 AI 生成图像越来越难以与真实照片区分；在一个生成器家族上表现准确的检测器，迁移到另一个家族时性能往往急剧下降。本文提出一种四分支注意力融合网络（FBFN），通过融合四种物理可解释的证据分支——CNN 语义分支、ViT 语义分支、基于误差水平分析（ELA）与隐写分析富模型（SRM）滤波器的噪声分支、以及基于 Lab 空间局部熵的颜色分支——并经过通道-空间双重注意力融合模块与全局池化，将图像分类为真实或生成。在 Stable Diffusion v1.4 子集上训练、并在 GenImage 协议下跨八个生成器评估，FBFN 在所有 Stable Diffusion 家族生成器上达到 99\% 以上的 ROC 曲线下面积（AUC），在该家族上与一个更强的近期黑盒检测器相当，同时保留了像素级可解释证据。消融实验表明，噪声分支在训练分布内近乎冗余，但具体负责跨生成器泛化，移除后在远距离生成器上 AUC 下降 5--7 个点。多生成器联合训练进一步弥合了大部分跨家族差距。我们讨论了该方法仍然薄弱之处——尤其是 GAN 生成图像——并将其可解释性定位为区别于更强位平面检测器的主要优势。
\end{abstract}

\noindent\textbf{关键词：}AI 生成图像检测；伪造检测；可解释深度学习；跨生成器泛化；注意力融合。

\section{引言}

大规模文本到图像模型现在可以低成本合成逼真图像，其输出越来越多地被用于虚假信息、身份欺诈和伪造证据。因此，检测这类图像已成为一个紧迫的实际问题，并催生了大量文献。核心难点不在于识别训练中见过的生成器生成的图像——这在很大程度上已经解决——而在于识别训练中未见过的新生成器生成的图像。在一个扩散模型上训练的检测器，面对来自不同家族的图像时往往失效，因为区分生成图像与真实图像的低级伪影在不同架构之间差异很大。

先前工作从两个主要方向攻克这一泛化难题。一条路线使用卷积或视觉 Transformer 骨干网络的通用语义特征 [1, 2]，泛化能力中等，但仍受限于训练分布的纹理统计。另一条路线利用手工设计或学习到的低级噪声线索——误差水平分析（ELA）[3]、隐写分析富模型（SRM）滤波器 [4]，以及最近的图像最低有效位平面 [5]——其假设是生成模型会在残差信号中留下特征性签名。这条路线的方法已被证明是迄今最强的跨生成器检测器，但它们通常是黑盒，只输出一个标量分数，不提供可供人类或下游系统检查的任何证据。

本文贡献了一种四分支架构，融合这两类证据，同时保持每个分支物理可解释。该网络我们称之为四分支注意力融合网络（FBFN），由 (i) CNN 分支和 (ii) ViT 分支捕获语义外观，(iii) 噪声分支捕获 ELA 和 SRM 残差伪影，(iv) 颜色分支捕获颜色空间和局部熵统计。四张像素级证据图通过通道-空间双重注意力融合模块组合，并池化为图像级分数。我们在 GenImage 协议下评估 FBFN——在 Stable Diffusion v1.4 上训练、跨八个生成器测试——并将其与三个代表性基线和一个新的更强检测器对比。我们进一步通过在训练分布内和分布外进行的留一法消融，剖析每个分支的贡献。

我们的结果说明了三点。第一，FBFN 在所有 Stable Diffusion 家族生成器上达到 99\% 以上的 AUC，与我们测试的最强黑盒检测器相当，而通用 RGB 基线的退化要大得多。第二，噪声分支——尽管在训练分布内近乎冗余——是负责大部分跨生成器鲁棒性的组件；移除它在最远距离生成器上使 AUC 下降 5--7 个点。第三，多生成器联合训练为扩散模型弥合了大部分剩余跨家族差距，但对 GAN 图像无帮助，该方法在此与若干基线一样接近随机。我们讨论了这些发现的含义，以及使 FBFN 区别于更强但不透明检测器的可解释性优势。

\section{相关工作}

\textbf{语义特征检测器。} 最早也最简单的方法，把生成图像检测当作使用标准图像骨干的二分类问题。GenImage 基准 [1] 建立了在单个生成器上训练的 ResNet-50 基线，并表明准确率跨生成器急剧下降。CNNDetection [2] 在训练时对 ResNet-50 加入模糊和 JPEG 损坏以诱导鲁棒性，但其零样本跨生成器准确率仍然有限。这些方法易于训练，但继承了骨干网络所学纹理统计的分布敏感性。

\textbf{低级伪影检测器。} 一条并行路线针对生成管线留下的残差信号。误差水平分析 [3]（可视化 JPEG 压缩不一致）和隐写分析富模型高通滤波器 [4]（最初为隐写术设计）都暴露了真实与生成图像之间不同的像素级证据。该路线最强的近期方法 LOTA [5] 提取图像的最低有效位平面，按梯度幅值选择图像块，并用 ResNet-50 分类；其报告的跨生成器准确率显著高于语义骨干。这些基于伪影的检测器泛化良好，但通常只产生一个标量决策，可解释性有限。

\textbf{零样本与基础模型检测器。} 第三条路线利用 CLIP [6] 等大型预训练视觉-语言模型的特征，假设生成图像落在自然图像流形之外。代表性方法 UnivFD [7] 仅在冻结的 CLIP 特征上训练一个线性头。这类检测器对某些生成器家族迁移良好，但对其他家族不敏感，反映了基础模型特征是为语义而非取证线索优化的这一事实。

FBFN 将前两条路线——语义证据与低级证据——结合在单一可解释架构中，与标量黑盒伪影检测器和通用骨干都不同：每个分支都产生一张可供人类检查的证据图。

\section{方法}

\subsection{任务定义}

我们将问题视为二分类图像分类。给定输入 RGB 图像 $x$，网络输出一个标量 logit，经过 sigmoid 后估计 $x$ 是 AI 生成的概率。模型在真实与生成图像对上以二值交叉熵监督训练。

\subsection{概述}

FBFN 的设计受人类视觉取证中``双过程认知''的启发：判断一张图像是否为生成，既依赖快速、直觉式的低级异常感知（如压缩块不连续、边缘跳变），也会在必要时调动更深层的语义推理（如物体结构是否合理）。认知科学将这一分工概括为系统一（快速、自动、低开销）与系统二（缓慢、审慎、高开销）[13]，近期也被引入 AI 系统计算分配的设计中 [14]。

FBFN 将这一划分映射到证据层面。四个分支被选择来覆盖互补且物理可解释的证据来源，每个分支产生一张空间尺寸相同的像素级证据图，随后是双重注意力融合模块与全局平均池化：

\begin{enumerate}
\item \textbf{CNN 分支}（系统二）—— 以 ResNet-34 [8] 为编码器的编码器-解码器网络（UNet++ [11]），捕获多尺度语义外观。
\item \textbf{ViT 分支}（系统二）—— Swin Transformer [9]（Tiny）搭配 SegFormer [10] MLP 解码器，捕获长程与全局语义结构。
\item \textbf{噪声分支}（系统一）—— 误差水平分析（JPEG 质量 90）与三个 5$\times$5 零和 SRM 高通滤波器拼接，捕获压缩与残差噪声伪影。
\item \textbf{颜色分支}（系统一）—— Lab 颜色空间特征加滑窗局部熵，后接浅层 CNN，捕获颜色统计异常。
\end{enumerate}

其中，系统一分支直接对低级残差信号建模，计算开销低、物理可解释；系统二分支对 RGB 语义内容建模，判别能力强但开销更高。这一划分并非先验假设，而是由消融结果倒推得出的（见 4.4 节）：系统一证据——尤其是噪声分支——是跨生成器泛化的主要来源，而系统二证据在同分布内提供强判别；二者的融合同时获得稳健性与可解释性。

\begin{figure}[h]
\centering
\includegraphics[width=0.98\textwidth]{\figdir/architecture.png}
\caption{所提出的四分支注意力融合网络（FBFN）总体结构。CNN 与 ViT 分支捕获语义外观，噪声（ELA + SRM）与颜色（Lab + 熵）分支捕获物理可解释的低级伪影；每个分支产生一张像素级证据图。四张图由通道-空间双重注意力模块融合，并全局池化为图像级真/假分数。}
\end{figure}

\subsection{证据分支}

延续上述双过程划分，系统一分支直接对低级残差信号建模，系统二分支对 RGB 语义内容建模；四者均将输入映射为单通道证据图。CNN 与 ViT 分支直接作用于 RGB 输入，产出语义证据。噪声分支接收 ELA 与 SRM 滤波响应作为辅助输入，并将其聚合为噪声证据。颜色分支接收熵图作为辅助输入，聚合颜色空间与熵证据。使用显式、命名的证据流——而非单一融合表示——正是使模型决策可检查之处：每张证据图都可被可视化，以显示网络依赖\textit{何处}以及\textit{何种}伪影。

\begin{figure}[h]
\centering
\includegraphics[width=0.9\textwidth]{\figdir/evidence_map.png}
\caption{一张真实图像（上）和一张生成图像（下）的分支证据图。每个分支输出单通道伪造概率图（暖色 = 伪造概率更高）；融合图综合四者。噪声与颜色分支突出了 RGB 输入中不可见的残差与颜色统计伪影，使决策可检查。}
\end{figure}

\subsection{双重注意力融合}

四张证据图拼接后通过一个融合模块，该模块包含两个协同的注意力机制。通道注意力路径采用全局平均池化、两层 MLP 和 softmax，对四种证据类型的重要性重新加权。空间注意力路径采用一个卷积后接 sigmoid，对每个空间位置的重要性重新加权。两个注意力输出合并后卷积为单张融合证据图。全局平均池化随后将融合图压缩为图像级 logit，另有一个注意力池化 logit 在训练时作为辅助信号。

\subsection{训练}

模型以组合二值交叉熵损失训练：

\begin{equation}
L = \mathrm{BCE}(\mathrm{final}) + 0.5\cdot\mathrm{BCE}(\mathrm{att}) + 0.5\cdot\sum \mathrm{BCE}(\mathrm{branch}),
\end{equation}

其中分支项对保留的分支求和。我们使用 AdamW（学习率 $10^{-4}$，权重衰减 $10^{-4}$）、余弦学习率调度、混合精度与水平翻转增强训练。输入图像缩放到 256$\times$256。完整模型含 54.68M 可训练参数。

\section{实验}

\subsection{设置}

\textbf{数据集与协议。} 我们遵循 GenImage 协议：模型仅在 Stable Diffusion v1.4 子集上训练，并在八个生成器（Stable Diffusion v1.4 与 v1.5、Wukong、ADM、GLIDE、BigGAN、Midjourney 与 VQDM）的验证集上评估。图像取自 TIGAS 数据集 [12]（GenImage [1] 的 Hugging Face 镜像），该数据集将 Stable Diffusion v1.5 子集拆成两部分；因此表格将 Stable Diffusion v1.5 报告为两行。训练集由 Stable Diffusion v1.4 的 6,000 张图像（3,000 真实、3,000 生成）组成；同分布验证集含 1,000 张图像；八个跨生成器验证子集各含 400--1,200 张图像（BigGAN 最少，为 200 真实与 200 生成）。

\textbf{指标。} 我们以 ROC 曲线下面积（AUC）为主指标，因为它与阈值无关，因而对生成器子集间不同的类别平衡不太敏感。准确率与平均精度（AP）均已计算，但为简洁起见不在表格中列出。

\textbf{基线。} 我们与在相同 Stable Diffusion v1.4 数据上训练的 ResNet-50、CNNDetection（零样本）和 UnivFD（零样本，CLIP 特征）对比。所有方法都在相同的测试图像上评估。

\subsection{同分布与跨生成器结果}

FBFN 在同分布 Stable Diffusion v1.4 验证集上达到 99.93\% 的 AUC。表~\ref{tab:1} 报告了所有方法的跨生成器 AUC。

\begin{table}[h]
\centering
\caption{跨生成器 AUC（\%），Stable Diffusion v1.4 训练。}
\label{tab:1}
\begin{tabular}{lccccc}
\toprule
生成器 & FBFN (sd14) & FBFN (联合) & ResNet-50 & CNNDetection & UnivFD \\
\midrule
SD v1.4 & 99.93 & --- & 97.70 & 56.96 & 73.54 \\
SD v1.5 (a) & 99.74 & --- & 98.28 & 60.00 & 72.59 \\
SD v1.5 (b) & 99.77 & --- & 97.24 & 59.89 & 74.61 \\
Wukong & 99.43 & --- & 96.55 & 53.24 & 80.55 \\
ADM & 66.04 & 83.87 & 37.34 & 69.00 & 94.30 \\
GLIDE & 68.42 & 94.48 & 52.86 & 70.65 & 86.39 \\
BigGAN & 54.23 & 48.31 & 36.62 & 88.11 & 99.00 \\
Midjourney & 76.03 & --- & 77.62 & 51.21 & 51.35 \\
VQDM & 77.40 & --- & 40.27 & 72.13 & 96.88 \\
\bottomrule
\end{tabular}
\end{table}

由此可得出三点观察。第一，FBFN 在 Stable Diffusion 家族和 Wukong 上近乎无损泛化（全部高于 99\% AUC），在该家族上大幅超越两个零样本基线，并超过 ResNet-50 基线。第二，在更远的扩散生成器（ADM、GLIDE、Midjourney、VQDM）上，FBFN 退化到 66--77\%，反映了语义分支未完全吸收的分布偏移。第三，在 BigGAN——一个 GAN 而非扩散模型——上，FBFN 接近随机（54.23\%），与其噪声和颜色证据针对扩散专属伪影的假设一致。值得注意的是，基线呈现出镜像模式：UnivFD 和 CNNDetection 在 BigGAN 上（99.00\% 与 88.11\%）远强于在 Stable Diffusion 家族上，证实这些方法捕获的是与 FBFN 不同种类的证据。

\subsection{联合训练}

为检验跨家族差距能否通过数据而非架构来缩小，我们在 Stable Diffusion v1.4、v1.5、VQDM 和 Midjourney 图像的组合池（共 9,000 张）上训练 FBFN，将 ADM、GLIDE 和 BigGAN 留作留出生成器。联合训练显著提升了留出扩散生成器——ADM 从 66.04\% 升至 83.87\%，GLIDE 从 68.42\% 升至 94.48\%——但没有提升 BigGAN（54.23\% $\to$ 48.31\%），后者仍停留在随机水平。这证实扩散模型的跨家族差距主要是数据覆盖差距，而 BigGAN 的失败反映了伪影类型上的更根本差异。

\subsection{消融研究}

我们进行了留一法消融，训练四个变体，各移除一个分支，同时保持相同协议（15 epoch，Stable Diffusion v1.4）。由于两个语义分支变体在远距离生成器上对训练种子敏感，我们各用不同随机初始化重跑一次；表~\ref{tab:2} 报告第二次运行，方差在下方量化。

\begin{table}[h]
\centering
\caption{消融 AUC（\%）。同分布（SD v1.4）与跨生成器。}
\label{tab:2}
\begin{tabular}{lccccc}
\toprule
配置 & SD v1.4 & ADM & GLIDE & BigGAN & VQDM \\
\midrule
Full & 99.93 & 66.04 & 68.42 & 54.23 & 77.40 \\
w/o CNN & 99.98 & 79.44 & 77.93 & 48.95 & 75.72 \\
w/o ViT & 99.40 & 57.67 & 57.05 & 52.71 & 70.48 \\
w/o Noise & 99.96 & 59.60 & 68.13 & 47.19 & 72.21 \\
w/o Color & 99.85 & 74.81 & 65.21 & 49.76 & 76.20 \\
\bottomrule
\end{tabular}
\end{table}

在训练分布内，所有变体都接近上限（AUC $\geq$ 99.4\%），因此无法从这些饱和分数中识别出某个占主导的分支。跨生成器结果则隔离出唯一稳健的贡献者。移除噪声分支使 AUC 在三个与训练分布距离最远的生成器上下降——ADM（$-6.44$）、BigGAN（$-7.04$）和 VQDM（$-5.19$）——而 Stable Diffusion 家族几乎不变。这表明 ELA/SRM 噪声证据充当了一个可跨生成器家族迁移的扩散伪影检测器，而非冗余的分布内特征。

其余分支没有显示出可靠的跨生成器贡献。CNN 与颜色分支在 ADM、GLIDE 和 VQDM 上仅改变几个点，且方向不一致（颜色分支在 ADM 上移除后甚至提升）。语义分支在远距离生成器上还表现出很大的运行间方差：两次运行中，w/o-CNN 的 GLIDE AUC 从 39.27\% 变为 77.93\%，w/o-ViT 的 GLIDE AUC 从 92.91\% 变为 57.05\%，摆幅约 35 个点且方向相反。因此，我们不将跨生成器贡献归因于 CNN 或 ViT 分支，并将其精确的跨生成器值视为不可靠。相比之下，噪声分支的退化在 ADM、BigGAN 和 VQDM 上符号与幅度都稳定，这正是使其成为跨生成器鲁棒性的可解释来源的原因。

\subsection{计算效率}

系统一/系统二的划分也体现为计算开销上的非对称分布。表~\ref{tab:3} 在 256$\times$256 输入下统计各分支参数量；表~\ref{tab:4} 报告在 NVIDIA RTX 5060 Laptop 上测得的前向延迟（50 次均值，GPU 预热后）。测量时前向仅传入预计算的 ELA 与熵图，以隔离网络本身的开销——ELA 编解码与滑窗熵属数据预处理，不计入网络前向。

\begin{table}[h]
\centering
\caption{各分支参数量（256$\times$256）。}
\label{tab:3}
\begin{tabular}{lcc}
\toprule
分支 & 参数量 (M) & 占比 \\
\midrule
噪声（系统一） & 0.04 & 0.07\% \\
颜色（系统一） & 0.41 & 0.75\% \\
CNN（系统二） & 26.08 & 47.7\% \\
ViT（系统二） & 28.15 & 51.5\% \\
\midrule
总计 & 54.68 & 100\% \\
\bottomrule
\end{tabular}
\end{table}

系统一分支合计 0.45M 参数（0.8\%），系统二分支合计 54.23M（99.2\%）；总 FLOPs 为 33.27 GFLOPs。

\begin{table}[h]
\centering
\caption{单图前向延迟（RTX 5060 Laptop，ms）。}
\label{tab:4}
\begin{tabular}{lc}
\toprule
配置 & 延迟 (ms) \\
\midrule
全模型（四分支+融合） & 30 \\
仅系统一（噪声+颜色） & 3 \\
仅系统二（CNN+ViT） & 27 \\
\bottomrule
\end{tabular}
\end{table}

系统一分支仅占总参数不到 1\%，单图延迟约 3 ms，而系统二分支承担了几乎全部计算开销，这直接印证了双过程类比中的成本假设：系统一证据几乎无成本，系统二证据昂贵。这一非对称性也自然引出一个级联推理方案——系统一先对全部图像快速筛选，仅当系统一无法给出高置信判断时才升级到全模型，期望平均延迟为 $T_1 + \rho\,(T - T_1)$，其中 $T_1 = 3$ ms 为系统一延迟，$T = 30$ ms 为全模型延迟，$\rho$ 为升级比例。表~\ref{tab:cascade} 在同分布验证集上报告了该级联的权衡。

\begin{table}[h]
\centering
\caption{同分布（SD v1.4）级联 AUC 与平均延迟权衡。}
\label{tab:cascade}
\begin{tabular}{lcc}
\toprule
升级比例 $\rho$ & 级联 AUC (\%) & 平均延迟 (ms) \\
\midrule
0.00（仅系统一） & 91.22 & 3 \\
0.38 & 97.56 & 13 \\
0.63 & 99.06 & 20 \\
0.85 & 99.66 & 26 \\
1.00（全模型） & 99.93 & 30 \\
\bottomrule
\end{tabular}
\end{table}

结果表明，系统一单独（噪声+颜色证据）在同分布上仅达 91.22\% AUC，且在跨生成器上普遍更弱（各生成器 36--69\% AUC）。级联要恢复到约 99\% AUC 需升级约 63\% 的样本，平均延迟约 20 ms——加速有限，且以明显精度损失为代价。这给出了一个与 4.4 节消融一致但更直接的结论：系统一证据虽然便宜，却不足以独立支撑高精度筛选；其价值在于与系统二证据融合以提供跨生成器鲁棒性，而非作为独立的快速通道。因此，除非训练一个更强的系统一头（例如为其补充轻量语义编码），否则级联在 FBFN 上的收益有限，我们将其留作未来工作。

\section{讨论}

结果对 FBFN 的优势来源与边界给出了一致的解释。在 Stable Diffusion 家族内及其间，FBFN 达到了我们所测试最强检测器的准确率，同时提供了标量检测器所缺乏的可解释证据结构。消融将噪声分支隔离为跨生成器鲁棒性的来源，联合训练则表明剩余的扩散家族差距是数据覆盖限制而非架构限制。

我们明确与 LOTA [5] 对比——这是一个近期的位平面检测器（ICCV 2025），在跨生成器 AUC 上强于 FBFN：LOTA 在 ADM 上达到 99.96\%、GLIDE 上 99.99\%、VQDM 上 100.00\%，而 FBFN 仅 66--77\%。我们不主张在这一轴上超越 LOTA。区别在于可解释性：LOTA 与其他伪影检测器一样，只输出标量分数，而 FBFN 的每个分支都产生一张像素级证据图，可检查是何种伪影、在哪个区域驱动了决策。对于决策必须可审计的应用——内容审核、取证分诊——即使原始跨生成器准确率较低，这种可审计性也是一项具体优势。

有两点局限值得强调。第一，该方法对 GAN 生成图像较弱，在 BigGAN 上退化到随机，且联合训练无法恢复；噪声与颜色证据似乎专属于扩散伪影。第二，我们的评估使用了 GenImage 的下采样子集（256$\times$256，6,000 张训练图像），这可能低估了完整基准的挑战，并限制了与在完整数据集上报告的数字的直接比较。尽管如此，本文所有比较都在相同数据上进行，因此相对发现不受影响。

\section{结论}

我们提出了 FBFN，一个用于 AI 生成图像检测的四分支注意力融合网络，其分支对应四种物理可解释的证据流。在 GenImage 协议下，FBFN 在 Stable Diffusion 家族上近乎无损泛化，且其噪声分支被消融证明具体负责跨生成器鲁棒性。联合训练弥合了大部分扩散家族差距，而 GAN 图像仍是一个未解决的失败模式。主要贡献不是跨生成器准确率的新最优——存在更强的黑盒检测器——而是一个在 Stable Diffusion 家族上保持竞争力、同时保留可解释证据的检测器。这种可解释性旨在支持标量分数不足时的可审计检测。

\section*{致谢}

本研究受 \textit{[基金来源与项目编号——待作者填写]} 资助。

\section*{利益冲突}

作者声明不存在利益冲突。

\section*{数据可用性声明}

本研究使用的 GenImage 基准可通过 TIGAS 镜像公开获取：\url{https://huggingface.co/datasets/H1merka/TIGAS_dataset}。训练模型与评估代码可向通讯作者合理索取。

\begin{thebibliography}{14}

\bibitem{1} Zhu, M., Chen, H., Yan, Q., et al.: ``GenImage: A million-scale benchmark for detecting AI-generated image'' \textit{Adv. Neural Inf. Process. Syst.}, 2023, 36.

\bibitem{2} Wang, S.-Y., Wang, O., Zhang, R., Owens, A., Efros, A.A.: ``CNN-generated images are surprisingly easy to spot... for now'' \textit{Proc. IEEE/CVF Conf. Comput. Vis. Pattern Recognit.}, 2020, pp. 8695--8704.

\bibitem{3} Krawetz, N.: ``A picture's worth: Digital image analysis and forensics'' \textit{Hacker Factor Blog}, 2007, \url{http://www.hackerfactor.com/blog/}, accessed August 2026.

\bibitem{4} Fridrich, J., Kodovsk\'{y}, J.: ``Rich models for steganalysis of digital images'' \textit{IEEE Trans. Inf. Forensics Secur.}, 2012, 7, (3), pp. 868--882.

\bibitem{5} Wang, H., Cheng, R., Zhang, Y., Han, C., Gui, J.: ``LOTA: Bit-planes guided AI-generated image detection'' \textit{Proc. IEEE/CVF Int. Conf. Comput. Vis.}, 2025.

\bibitem{6} Radford, A., Kim, J.W., Hallacy, C., et al.: ``Learning transferable visual models from natural language supervision'' \textit{Proc. Int. Conf. Mach. Learn.}, 2021, 139, pp. 8748--8763.

\bibitem{7} Ojha, U., Li, Y., Lee, Y.J.: ``Towards universal fake image detectors that generalize across generative models'' \textit{Proc. IEEE/CVF Conf. Comput. Vis. Pattern Recognit.}, 2023, pp. 24480--24489.

\bibitem{8} He, K., Zhang, X., Ren, S., Sun, J.: ``Deep residual learning for image recognition'' \textit{Proc. IEEE/CVF Conf. Comput. Vis. Pattern Recognit.}, 2016, pp. 770--778.

\bibitem{9} Liu, Z., Lin, Y., Cao, Y., et al.: ``Swin transformer: Hierarchical vision transformer using shifted windows'' \textit{Proc. IEEE/CVF Int. Conf. Comput. Vis.}, 2021, pp. 10012--10022.

\bibitem{10} Xie, E., Wang, W., Yu, Z., Anandkumar, A., Alvarez, J.M., Luo, P.: ``SegFormer: Simple and efficient design for semantic segmentation with transformers'' \textit{Adv. Neural Inf. Process. Syst.}, 2021, 34, pp. 12077--12090.

\bibitem{11} Zhou, Z., Rahman Siddiquee, M.M., Tajbakhsh, N., Liang, J.: ``UNet++: A nested U-Net architecture for medical image segmentation'' \textit{Deep Learn. Med. Image Anal. Multimodal Learn. Clin. Decis. Support}, 2018, pp. 3--11.

\bibitem{12} H1merka: ``TIGAS dataset'' \textit{Hugging Face}, \url{https://huggingface.co/datasets/H1merka/TIGAS_dataset}, accessed August 2026.

\bibitem{13} Kahneman, D.: ``Thinking, fast and slow'' \textit{Farrar, Straus and Giroux}, 2011.

\bibitem{14} Weston, J., Sukhbaatar, S.: ``System 2 attention (is something you might need too)'' \textit{arXiv preprint arXiv:2311.11829}, 2023.

\end{thebibliography}

\end{document}
